\section*{Bab \ref{ch:b3:glaciology} --- Glasiologi}

\begin{solution}{exo:b3:glaciology:1}
$\tau_b = \rho gH\sin\alpha = 917\times9.81\times250\times\sin\ang{3} =
\qty{118}{kPa}$, sekitar \qty{1.2}{bar}: orde tegangan yang
dialami gletser sungguhan.
\end{solution}

\begin{solution}{exo:b3:glaciology:2}
(a)~$\qty{2.9e15}{m^3}\times917 = \qty{2.66e18}{kg} = \qty{2.66e6}{Gt}$;
dibagi \qty{360}{Gt/mm}, \qty{7390}{mm}, \qty{7.4}{m}, nilai yang dikutip.
(b)~$\qty{3.0e16}{m^3}\times917/360 = \qty{76}{m}$, dibandingkan dengan \qty{58}{m}.
Es di Greenland sebagian besar bertumpu pada batuan dasar di atas permukaan laut, sehingga seluruhnya
akan menambah air ke samudra. Sebagian besar es Antarktika bertumpu pada dasar di bawah
permukaan laut: es itu sudah memindahkan air laut, dan mencairkannya hanya menggantikan
air yang dipindahkannya; rak es terapung tidak menambah apa pun sama sekali.
\end{solution}

\begin{solution}{exo:b3:glaciology:3}
Penambahan $\num{6e6}\times1.2 = \qty{7.2e6}{m^3}$, kehilangan $\num{4e6}\times2.1 =
\qty{8.4e6}{m^3}$ air: bersih \qty{-1.2e6}{m^3}, rata-rata neraca spesifik
sebesar $-1.2\times10^6/10^7 = \qty{-0.12}{m}$ setara air.
Rasio luas \omterm{def:b3:glaciology:mass-balance}{zona akumulasi} adalah $6/10 = 0.6$. Gletser ini kehilangan massa
dan menyusut.
\end{solution}

\begin{solution}{exo:b3:glaciology:4}
$1 - 820/917 = \num{0.106}$ dan $1 - 830/917 = \num{0.095}$: sekitar sepersepuluh
dari volume masih berupa ruang pori ketika pori-pori menutup.
\end{solution}

\begin{solution}{exo:b3:glaciology:5}
$\tau_b = 917\times9.81\times300\times\sin\ang{3} = \qty{141.2}{kPa}$. Menurut
teorema lempengan, $u_s = \tfrac{2A}{4}\tau_b^3H = 0.5\times\num{1.43e-16}
\times(\num{1.412e5})^3\times300 = \qty{60}{m/yr}$. Bila tebal dua kali lipat: $u_s
\propto H^{n+1} = H^4$, enam belas kali lebih cepat. Jika kemiringan setengahnya: $u_s \propto
(\sin\alpha)^3$, delapan kali lebih lambat.
\end{solution}

\begin{solution}{exo:b3:glaciology:6}
Pada \qty{-25}{\celsius} (\qty{248.15}{K}), $A = \num{3.61e-13}\exp(-60000/
(8.314\times248.15)) = \qty{8.5e-26}{Pa^{-3}.s^{-1}}$, atau
\qty{2.7e-18}{Pa^{-3}.yr^{-1}}. Pada \qty{-10}{\celsius}, $\num{3.61e-13}
\exp(-27.42) = \qty{4.4e-25}{Pa^{-3}.s^{-1}} = \qty{1.4e-17}{Pa^{-3}.yr^{-1}}$.
Pada \qty{0}{\celsius} (cabang hangat) \qty{1.4e-16}{Pa^{-3}.yr^{-1}}. Di bawah
tegangan yang sama, kecepatannya berbanding sebagai rasio \omterm{def:b3:glaciology:glen}{faktor laju}:
$\num{1.43e-16}/\num{2.67e-18} = 54$.
\end{solution}

\begin{solution}{exo:b3:glaciology:7}
$h = \sqrt{2\times10^5\times10^6/(917\times9.81)} = \qty{4.7}{km}$, mendekati
\qty{4.9}{km} di Antarktika: model plastis memperoleh ketebalan sebuah
lembaran es hanya dari lebarnya.
\end{solution}

\begin{solution}{exo:b3:glaciology:8}
$\tau = 200/5 = \num{40}\text{ tahun}$. Panjangnya akan terus menyesuaikan diri selama
beberapa dekade, beberapa kali waktu respons, setelah pemanasan berhenti: gletser saat ini
belum menyesuaikan diri dengan iklim saat ini.
\end{solution}

\begin{solution}{exo:b3:glaciology:9}
(a)~$h = \sqrt{2\times2.03\times20\times\num{8.64e6}/(917\times\num{3.34e5})}
= \qty{1.51}{m}$. Karena $h^2 \propto t$, \qty{3}{m} membutuhkan $100\times
(3/1.51)^2 = 393$ hari, lebih dari satu musim dingin. (b)~Es yang jauh lebih tebal daripada
\qty{3}{m} tidak dapat tumbuh dengan membeku dalam satu musim dingin: es yang
telah bertahan melewati beberapa musim, atau es yang ditumpuk angin dan arus di punggungan.
\end{solution}

\begin{solution}{exo:b3:glaciology:10}
(a)~$H = (1025/917)\times800 = \qty{894}{m}$. (b)~Pada \qty{900}{m}
tebal \omterm{def:b3:glaciology:grounding-line}{garis kandas} adalah \qty{1006}{m}, $9/8$ dari sebelumnya; fluks
naik sebesar $(9/8)^5 = 1.8$. (c)~Ketika garis mundur menuruni dasar yang makin dalam
ke arah daratan, es di garis itu lebih tebal, debit lebih besar, menipiskan es
di belakangnya, lalu garis itu mundur lebih jauh: umpan balik positif. Paparan es
yang menopang es bertumpu di dasar mengurangi fluks di \omterm{def:b3:glaciology:grounding-line}{garis kandas} dan dapat
menahannya tetap di tempat walaupun pada lereng yang menurun.
\end{solution}

\begin{solution}{exo:b3:glaciology:11}
Dengan $h(r) = c\sqrt{R - r}$, $c = \sqrt{2\tau_0/(\rho g)}$, $V =
\int_0^R2\pi rc\sqrt{R - r}\,\mathrm dr$. Substitusi $s = R - r$,
\[
\int_0^R(R - s)\sqrt s\,\mathrm ds = \tfrac23R^{5/2} - \tfrac25R^{5/2} =
\tfrac4{15}R^{5/2},
\]
sehingga $V = \frac{8\pi}{15}cR^{5/2}$. Menggandakan $R$
mengalikan $V$, dan \omterm{def:b3:glaciology:sle}{setara permukaan laut}, dengan $2^{5/2} = 5.7$.
\end{solution}

\begin{solution}{exo:b3:glaciology:12}
(a)~Maret: $14.13/16.34 - 1 = \qty{-13.5}{\percent}$; September: $4.81/7.05
- 1 = \qty{-31.8}{\percent}$. (b)~Pada Maret tepi es sebagian besar ditahan oleh
pesisir yang mengelilingi Samudra Arktik, dan pembekuan ulang setiap musim dingin
mengisi cekungan kembali; es yang lebih tipis dalam beberapa tahun terakhir mencair lebih
tuntas setiap musim panas. Umpan balik es--albedo juga bekerja pada musim panas: perairan terbuka
menyerap sekitar \qty{94}{\percent} sinar Matahari, es memantulkan
\qtyrange{50}{70}{\percent}, sehingga setiap kehilangan es musim panas memanaskan air dan
mencair lebih banyak.
\end{solution}

\begin{solution}{pb:b3:glaciology:1}
\textbf{1.} Kolom di atas dasar, dengan berat $\rho gH$ per satuan luas, didorong
ke arah bawah lereng oleh komponen $\rho gH\sin\alpha$, diseimbangkan oleh
tegangan geser pada alasnya; untuk lereng kecil $\sin\alpha \approx \alpha$ (teorema
lempengan dengan kemiringan permukaan).

\textbf{2.} $917\times9.81\times580\times0.0190 = \qty{99.1}{kPa}$;
$917\times9.81\times470\times0.0240 = \qty{101.5}{kPa}$; $917\times9.81
\times335\times0.0330 = \qty{99.4}{kPa}$.

\textbf{3.} Tegangan tetap sama dalam \qty{2}{\percent} sementara
ketebalannya hampir separuh: tudung itu mendekati plastis, dengan $\tau_0 \approx
\qty{100}{kPa}$.

\textbf{4.} $2\times10^5/(917\times9.81) = \qty{22.23}{m}$.

\textbf{5.} $h = \sqrt{22.23\,(R - r)}$: \qty{577}{m}, \qty{472}{m} dan
\qty{333}{m} pada 15, 10, dan \qty{5}{km} dari tepi, dalam beberapa meter
dari data pengukuran.

\textbf{6.} $\sqrt{22.23\times20000} = \qty{667}{m}$.

\textbf{7.} Seperti pada \cref{exo:b3:glaciology:11}: $V = \int_0^R2\pi r
\sqrt{22.23(R - r)}\,\mathrm dr = \frac{8\pi}{15}\sqrt{2\tau_0/(\rho
g)}\,R^{5/2}$.

\textbf{8.} $\frac{8\pi}{15}\times4.715\times(\num{2e4})^{5/2} =
\qty{4.47e11}{m^3} = \qty{447}{km^3}$.

\textbf{9.} $\pi\times20^2 = \qty{1257}{km^2}$; ketebalan rata-rata
$447/1257 = \qty{0.356}{km}$, \qty{356}{m}, sedikit lebih dari separuh
ketebalan di pusatnya.

\textbf{10.} $2^{5/2} = 5.7$ kali.

\textbf{11.} $A = \num{3.61e-13}\exp(-60000/(8.314\times263.15)) =
\qty{4.45e-25}{Pa^{-3}.s^{-1}}$, dikali \qty{3.156e7}{s}:
\qty{1.40e-17}{Pa^{-3}.yr^{-1}}.

\textbf{12.} Pada \qty{10}{km}, $H = \qty{471.5}{m}$, kemiringan $471.5/20000 =
\num{0.0236}$, $\tau_b = \qty{100}{kPa}$; $u_s = \tfrac{2A}{4}\tau_b^3H =
0.5\times\num{1.40e-17}\times10^{15}\times471.5 = \qty{3.3}{m/yr}$.

\textbf{13.} $q = \tfrac{2A}{5}\tau_b^3H^2 = 0.4\times\num{1.40e-17}
\times10^{15}\times471.5^2 = \qty{1250}{m^2/yr}$; kecepatan rata-rata $q/H =
\qty{2.6}{m/yr}$, empat per lima dari kecepatan permukaan.

\textbf{14.} $2\pi rq = \pi r^2\bar b$, sehingga $\bar b = 2q/r = 2\times1250/10^4
= \qty{0.25}{m}$ es per tahun pada \qty{10}{km} bagian dalam.

\textbf{15.} $h = 450$ bila $R - r = 450^2/22.23 = \qty{9.11}{km}$, pada
$r = \qty{10.9}{km}$; \omterm{def:b3:glaciology:mass-balance}{zona akumulasi} adalah cakram di dalamnya, $(10.9/
20)^2 = 0.30$ dari tudung.

\textbf{16.} $A$ naik sebesar \num{10.2}, dan begitu pula kecepatan di bawah
tegangan yang sama. Es akan mengalir keluar lebih cepat daripada yang dipasok neraca massa: tudung
akan menipis dan melebar hingga tegangan yang lebih rendah kembali sesuai dengan
neraca massa; tudung plastis yang lebih lunak memiliki tegangan luluh yang lebih kecil dan
profil yang lebih datar.

\textbf{17.} $\num{4.47e11}\times917 = \qty{4.10e14}{kg} = \qty{410}{Gt}$.

\textbf{18.} $410/360 = \qty{1.14}{mm}$.

\textbf{19.} $356/0.5 = 711$ tahun.

\textbf{20.} Ketika permukaan turun, ia berada di udara yang lebih hangat dan mencair lebih cepat,
umpan balik elevasi--neraca massa; dan ketika menipis, tegangan dan
aliran juga ikut berubah. Penipisan pun makin cepat.

\textbf{21.} $410/243 = \num{1.7}\text{ tahun}$: seluruh tudung kurang dari dua
tahun kehilangan massa Greenland pada 2010--2019.

\textbf{22.} Tudung es Arkel memiliki \omterm{def:b3:glaciology:sle}{setara permukaan laut} sebesar
\textbf{\qty[text-series-to-math=true]{1.1}{mm}} (\qty{1.14}{mm}).
\end{solution}
